\documentclass[11pt]{amsart}
\usepackage[margin=1.1in]{geometry}
\usepackage{amssymb,amsmath,amsthm,booktabs,array,xcolor,url}
\usepackage[hidelinks]{hyperref}
\newcommand{\TODO}[1]{\textcolor{red}{\textbf{[TODO: #1]}}}
\newtheorem{theorem}{Theorem}
\newtheorem{lemma}[theorem]{Lemma}
\newtheorem{proposition}[theorem]{Proposition}
\theoremstyle{definition}
\newtheorem{definition}[theorem]{Definition}
\theoremstyle{remark}
\newtheorem{remark}[theorem]{Remark}
\newcommand{\R}{R}
\newcommand{\lexle}{\preceq_{\mathrm{lex}}}

\title[The cycle Ramsey numbers $R(C_4,C_7,C_7)$ and $R(C_6,C_7,C_7)$]{The three-colour cycle Ramsey numbers\\ $R(C_4,C_7,C_7)$ and $R(C_6,C_7,C_7)$}
\author{Larsen Cundric}
\address{\TODO{affiliation and postal address}}
\email{\TODO{e-mail}}
\date{Draft of \today}
\subjclass[2020]{05C55, 05D10, 68V15}
\keywords{Ramsey numbers, cycles, SAT solving, cube-and-conquer, proof certificates}

\begin{document}

\begin{abstract}
For graphs $G_1,G_2,G_3$, the Ramsey number $R(G_1,G_2,G_3)$ is the least $N$ such that every $3$-colouring of the edges of
$K_N$ contains a copy of $G_i$ in colour $i$ for some $i$. In the table of three-colour cycle Ramsey numbers
$R(C_m,C_n,C_k)$ with $m,n,k\le 7$ in Radziszowski's dynamic survey, the only entries without a value are
$R(C_m,C_7,C_7)$ for $m=3,4,5,6$. We show that $R(C_4,C_7,C_7)=15$ and \TODO{pending: $R(C_6,C_7,C_7)=21$}.
The lower bounds come from explicit colourings. The upper bounds come from a SAT-based cube-and-conquer
computation in which forbidden-cycle clauses are generated lazily. Every part of it is certified: each clause the solver uses
is independently checked to be a genuine forbidden-cycle, colouring, or symmetry-breaking clause; every subproblem has an
unsatisfiability proof checked by the formally verified checker \texttt{cake\_lpr}; and the cube split is covered by a separate
checked certificate. The result for $R(C_4,C_7,C_7)$ was also obtained a second time with an independent encoding, symmetry-breaking
scheme, lazy solver, and proof checker (DRAT with \texttt{drat-trim}) \TODO{confirm when leg B finishes}. We also give
colourings showing $R(C_3,C_7,C_7)\ge 25$ and $R(C_5,C_7,C_7)\ge 25$.
\end{abstract}

\maketitle

\section{Introduction}

All graphs are finite and simple, and $C_k$ is the cycle of length $k$. For graphs $G_1,\dots,G_r$ the Ramsey number
$R(G_1,\dots,G_r)$ is the least $N$ such that every colouring of $E(K_N)$ with colours $1,\dots,r$ has, for some $i$, a copy of
$G_i$ all of whose edges have colour $i$. A colouring of $K_{N-1}$ with no such copy is called \emph{good}; it shows $R>N-1$.

Cycle Ramsey numbers in three colours have a long history; Radziszowski's dynamic survey~\cite{Rad} (revision~18, April 2026)
collects all known values in its Section~6.3.1 and Table~XII. Erd\H{o}s, Faudree, Rousseau and Schelp~\cite{EFRS1} determined
$R(C_m,C_n,C_k)$ for large~$m$ (formulas quoted in~\cite[6.3.1(a)]{Rad}). For three equal odd cycles, the conjecture attributed to
Bondy and Erd\H{o}s that $R(C_n,C_n,C_n)=4n-3$ for odd $n\ge5$ (see~\cite[6.3.1]{Rad}) was proved for large~$n$ by
Kohayakawa, Simonovits and Skokan~\cite{KoSS2}, and $R(C_7,C_7,C_7)=25$ was proved without computer support by Faudree, Schelten and
Schiermeyer~\cite{FSS3}. Most of the small off-diagonal values in Table~XII of~\cite{Rad} were obtained by computer, mainly by Sun,
Yang, Wang, Li and Xu~\cite{SunYWLX} and Sun, Yang, Lin and Zheng~\cite{SunYLZ2}, and independently by Tse~\cite{Tse3}. More
recently, Lidick\'y and Pfender~\cite{LidP} used flag algebras, and Wesley~\cite{We3} used SAT solvers to prove
$R(C_3,C_6,C_6)=R(C_5,C_6,C_6)=15$.

In Table~XII of~\cite{Rad} (all $m\le n\le k\le 7$, plus $m=n=k=8$), exactly four entries are blank:
$R(C_3,C_7,C_7)$, $R(C_4,C_7,C_7)$, $R(C_5,C_7,C_7)$ and $R(C_6,C_7,C_7)$. According to the abstract of~\cite{SunYWLX},
that work treats $R(C_{m_0},C_{m_1},C_{m_2})$ with $m_2\le m_1<m_0\le 7$, so the case where the two longer cycles have equal length~$7$ falls outside it.
We searched arXiv, Google Scholar, MathOverflow and Chinese-language sources (as of October~2026) and found no claimed value for any
of the four entries; we cannot rule out unindexed Chinese-language work.

\begin{theorem}\label{thm:main}
$R(C_4,C_7,C_7)=15$.
\end{theorem}

\begin{theorem}\label{thm:c6}
\TODO{PENDING (computation running): $R(C_6,C_7,C_7)=21$.} At present we prove $R(C_6,C_7,C_7)\ge 21$.
\end{theorem}

\begin{proposition}\label{prop:odd}
$R(C_3,C_7,C_7)\ge 25$ and $R(C_5,C_7,C_7)\ge 25$.
\end{proposition}

Here and throughout, the cycle lengths are assigned to colours $0,1,2$ in the order written: a colouring is good for
$(C_m,C_7,C_7)$ if colour~$0$ has no $C_m$ and colours $1$ and $2$ have no $C_7$.

The upper bound in Theorem~\ref{thm:main} is a finite but large computation; the full CNF encoding of ``$K_{15}$ has a good
colouring'' has about $4.6$ million clauses, and for $K_{21}$ the $C_7$-clauses alone number about $8.4\times10^7$. Our method
(Section~\ref{sec:sat}) avoids generating most of these clauses by adding forbidden-cycle clauses lazily, splits the problem with
cube-and-conquer~\cite{HKWB}, and certifies every step (Section~\ref{sec:cert}).

\section{Lower bounds}\label{sec:lower}

All colourings below are given as explicit adjacency matrices in the accompanying artifact (Section~\ref{sec:repro}). Each was
checked with two independent programs: one using a depth-first search for monochromatic $k$-cycles, the other using the
\texttt{simple\_cycles} routine of NetworkX~3.7 with a length bound. The two checkers were cross-validated on 2730 random colourings
(0 disagreements), including brute-force enumeration on small instances.

\begin{proposition}\label{prop:lb4}
There is a good $3$-colouring of $K_{14}$ for $(C_4,C_7,C_7)$. Hence $R(C_4,C_7,C_7)\ge15$.
\end{proposition}
\begin{proof}
Let $V=A\cup B\cup\{s_1,s_2\}$ with $|A|=|B|=6$. Colour the edges inside $A$ and inside $B$ with colour~$1$; colour $s_1a$ ($a\in A$)
and $s_2b$ ($b\in B$) with colour~$0$; colour all other edges ($A$--$B$, $s_1$--$B$, $s_2$--$A$, $s_1s_2$) with colour~$2$.
Colour~$0$ is two disjoint stars $K_{1,6}$, which are acyclic. Colour~$1$ is two disjoint copies of $K_6$, so it has no cycle on $7$
vertices. Colour~$2$ is bipartite with parts $A\cup\{s_1\}$ and $B\cup\{s_2\}$, so it has no odd cycle.
\end{proof}

\begin{proposition}\label{prop:lb6}
There is a good $3$-colouring of $K_{20}$ for $(C_6,C_7,C_7)$. Hence $R(C_6,C_7,C_7)\ge21$.
\end{proposition}
\begin{proof}
Partition $V$ into $P_0,P_1,P_2,P_3$ of size $5$. Edges inside a part get colour $0$ (four disjoint $K_5$; no cycle of length
$6$). Edges between $P_0\cup P_1$ and $P_2\cup P_3$ get colour~$1$ (a $K_{10,10}$, bipartite). Edges between $P_0$ and $P_1$ and between $P_2$
and $P_3$ get colour $2$ (two disjoint copies of $K_{5,5}$, bipartite). Bipartite graphs contain no $C_7$.
\end{proof}

\begin{proof}[Proof of Proposition~\ref{prop:odd}]
Partition $V(K_{24})$ into $P_0,\dots,P_3$ of size $6$. Inside parts use colour~$1$ (four disjoint $K_6$, no $C_7$);
$P_0$--$P_1$ and $P_2$--$P_3$ use colour~$2$ (bipartite); $(P_0\cup P_1)$--$(P_2\cup P_3)$ uses colour $0$ (a $K_{12,12}$, which has
no odd cycle, in particular no $C_3$ and no $C_5$).
\end{proof}

\begin{remark}
These are block constructions of the usual kind, and they agree with the values suggested by the neighbouring entries of Table~XII
of~\cite{Rad}: for example $R(C_3,C_5,C_7)=R(C_5,C_5,C_7)=R(C_7,C_7,C_7)=25$ and $R(C_6,C_6,C_7)=15$. An exhaustive search over
circulant colourings found no good circulant colouring of $K_{15}$ for $(C_4,C_7,C_7)$, of $K_{21}$ for $(C_6,C_7,C_7)$, or of
$K_{25}$ for $(C_3,C_7,C_7)$ or $(C_5,C_7,C_7)$.
\end{remark}

\section{SAT encodings}\label{sec:sat}

Fix $n$ and lengths $L=(\ell_0,\ell_1,\ell_2)$ (here $\ell_0=m$, $\ell_1=\ell_2=7$). Write $E=\binom{[n]}{2}$, ordered
lexicographically.

\subsection{Encoding A (one-hot)}
Variables $x_{e,c}$ for $e\in E$, $c\in\{0,1,2\}$, meaning ``edge $e$ has colour $c$''. Clauses:
\begin{itemize}
\item at-least-one: $x_{e,0}\vee x_{e,1}\vee x_{e,2}$ for every $e$;
\item at-most-one: $\neg x_{e,c}\vee\neg x_{e,c'}$ for $c<c'$;
\item forbidden cycles: for each colour $c$ and each $\ell_c$-cycle $Z$ of $K_n$, the clause $\bigvee_{e\in Z}\neg x_{e,c}$.
\end{itemize}
In a purely logical sense the at-most-one clauses are redundant (any model can be reduced to one colour per edge), but they speed up
solving considerably: on $K_{12}$ for $(C_4,C_4,C_7)$, Kissat takes 143\,s with them and more than 20 minutes without them. Cycles are
enumerated once each by a depth-first search from the smallest vertex, with the second vertex smaller than the last.

\subsection{Encoding B (independent)}
Variables $r_e, g_e$ for $e\in E$; colour $0$ is $r_e$, colour $1$ is $g_e$, colour $2$ is $\neg r_e\wedge\neg g_e$; clause
$\neg r_e\vee\neg g_e$. Forbidden cycles: $\bigvee_{e\in Z}\neg r_e$ for colour $0$, $\bigvee_{e\in Z}\neg g_e$ for colour $1$, and
$\bigvee_{e\in Z}(r_e\vee g_e)$ for colour $2$. Encoding~B was written separately from Encoding~A. Its generator enumerates cycles by
$k$-subsets and permutations, and its lazy solver uses its own depth-first search.

\subsection{Symmetry breaking}\label{sec:sb}
The symmetric group $S_n$ acts on colourings by permuting vertices, and good colourings are mapped to good colourings. Fix a total
order on the variables of the encoding, and for an assignment $\alpha$ write $\alpha$ as a $0/1$ vector in that order. For a vertex
permutation $\sigma$ let $\sigma\alpha$ be the assignment of the permuted colouring, i.e.\ $(\sigma\alpha)(v_{\{i,j\},c})=
\alpha(v_{\{\sigma^{-1}i,\sigma^{-1}j\},c})$.

\begin{lemma}[lex-leader, cf.~\cite{CGLR}]\label{lem:lexleader}
Let $T\subseteq S_n$ be any set of permutations and let $\Phi_T$ be a CNF over the original variables and fresh auxiliary variables such
that for every assignment $\alpha$ of the original variables with $\alpha\lexle\sigma\alpha$ for all $\sigma\in T$, some extension of
$\alpha$ to the auxiliary variables satisfies $\Phi_T$. If a good colouring of $K_n$ exists, then $F\wedge\Phi_T$ is satisfiable,
where $F$ is the full encoding (A or~B).
\end{lemma}
\begin{proof}
Let $\chi$ be good and let $\alpha^*$ be the lexicographically least assignment in the (finite) orbit $\{\sigma\alpha_\chi:\sigma\in
S_n\}$. Every element of the orbit encodes a good colouring and so satisfies $F$. For $\sigma\in T$, $\sigma\alpha^*$ also lies in the orbit, so
$\alpha^*\lexle\sigma\alpha^*$. By hypothesis $\alpha^*$ extends to a model of $\Phi_T$, and $F$ does not mention the auxiliary variables.
\end{proof}

We use $T=\{(i\ i{+}1):0\le i<n-1\}$ (adjacent transpositions). We encode $\alpha\lexle\sigma\alpha$ in the standard way: let
$X_1,\dots,X_p$ be the variables moved by $\sigma$ (in the variable order) and $Y_t$ the variable whose value $\sigma\alpha$ puts at
position $X_t$. With auxiliary variables $q_1,\dots,q_{p-1}$ (``equal so far''), the clauses are
\[
q_{t-1}\rightarrow(X_t\rightarrow Y_t),\qquad q_{t-1}\wedge(X_t\leftrightarrow Y_t)\rightarrow q_t\qquad(q_0:=\top).
\]
If $\alpha\lexle\sigma\alpha$, then setting $q_t$ true exactly when $X_s=Y_s$ for all $s\le t$ satisfies these clauses, so
Lemma~\ref{lem:lexleader} applies. Encodings A and~B use different variable sets and orders (A: $x_{e,0},x_{e,1},x_{e,2}$ per edge;
B: $r_e,g_e$ per edge), so they lead to different lexicographic orders and different symmetry-breaking predicates. As an additional
check of the implementation, we tested Encoding~B's predicates by brute force on $K_5$ and $K_6$: for random colourings, the
lex-least member of the full $S_n$-orbit always satisfied the predicates, while the lex-greatest member always violated them.

\subsection{Lazy cycle clauses}
For $\ell=7$ almost all clauses are cycle clauses: $K_{15}$ has $2{,}316{,}600$ seven-cycles, and $K_{21}$ has about
$4.19\times10^7$. We therefore start from a \emph{base formula}~$B$, which contains all colouring clauses, the symmetry-breaking
predicates, and the cycle clauses of length at most~$5$, and add longer cycle clauses only when needed:
\begin{itemize}
\item In leg~A we use an external propagator for CaDiCaL~\cite{CaDiCaL2} through the IPASIR-UP interface~\cite{IPASIRUP}. It tracks
the partial colouring and, whenever the edges currently set to colour~$c$ close an $\ell_c$-cycle, adds that cycle's clause (which is
then falsified) as a conflict clause. It also checks every candidate model.
\item In leg~B we use a counterexample-guided loop on top of CaDiCaL~1.9.5 (through PySAT). We solve; if the model contains forbidden
monochromatic cycles, we add up to $200$ of their clauses per colour and solve again; otherwise the model is a good colouring.
\end{itemize}
Every clause added in this way is a forbidden-cycle clause of the full encoding, so it is implied by the encoding's semantics, and
the procedure is sound and complete. The added clauses are not trusted. They are written to a file and re-checked (Section~\ref{sec:cert}).

\subsection{Cube-and-conquer}
We split the search space with the look-ahead cuber \texttt{march\_cu}~\cite{HKWB} at depth~$10$, run on $B$ together with the first
$10^5$ cycle clauses from a short lazy warm-up run. The split itself is not trusted: we only use the resulting list of cubes
$Q_1,\dots,Q_s$, i.e.\ conjunctions of literals. For $K_{15}$, leg~A has $s=988$ cubes and leg~B has $s=1022$.

\emph{Adaptive re-splitting.} A few cubes are far harder than the rest; for example, one $K_{21}$ cube took $32{,}392$\,s, while most
take seconds. A cube (more generally, a node of the cube tree) is therefore split again when its lazy solve exceeds $1200$\,s, or when
its offline proof run or proof check exceeds $2400$\,s. To split node $P$ we run \texttt{march\_cu} (depth~$6$) on $B$, together with up
to $10^5$ of the cycle clauses the lazy solver had found under $P$ and the literals of $P$ as units. The result is a list of subcubes
$D_1,\dots,D_k$. The children of $P$ are $P\wedge D_j$, and the split is accepted only with a checked coverage certificate (see (C2$'$)
below).

\section{Certification}\label{sec:cert}

For each leg the certificate of ``no good colouring of $K_n$'' consists of the following.
\begin{enumerate}
\item[(C1)] \emph{Base legitimacy.} A separate checker reads $B$ clause by clause and accepts it only if every clause is a colouring
clause (at-least-one/at-most-one, or exclusion for~B), a genuine forbidden-cycle clause (its literals form a single cycle of the
forbidden length, in the right colour), or a symmetry-breaking clause regenerated from scratch.
\item[(C2)] \emph{Coverage.} The formula $\bigwedge_i\neg Q_i$ is unsatisfiable, so every assignment satisfies some cube.
\item[(C2$'$)] \emph{Split coverage.} For each split of a node $P$ into $P\wedge D_1,\dots,P\wedge D_k$, the formula
$P\wedge\bigwedge_j\neg D_j$ ($P$ as unit clauses) is unsatisfiable. As a fallback, $B\wedge P\wedge\bigwedge_j\neg D_j$ is unsatisfiable.
\item[(C3)] \emph{Cubes.} For each leaf $Q_i$ of the cube tree, let $E_i$ be the clauses added by the lazy solver while working under the assumptions $Q_i$.
The checker verifies that every clause of $E_i$ is a genuine forbidden-cycle clause. The formula $F_i=B\wedge E_i\wedge Q_i$ (with $Q_i$
as unit clauses) is then solved offline from scratch with proof logging, and the proof is checked.
\end{enumerate}

\begin{theorem}\label{thm:cert}
If (C1)--(C3) hold, then $K_n$ has no good colouring.
\end{theorem}
\begin{proof}
Suppose $\chi$ is good, and let $\alpha^*$ be the lex-least assignment in its orbit, extended to the auxiliary variables as in
Section~\ref{sec:sb}. By (C1) and Lemma~\ref{lem:lexleader}, $\alpha^*$ satisfies $B$. Each clause of $E_i$ forbids a monochromatic
$\ell_c$-cycle in colour~$c$, and $\alpha^*$ encodes a good colouring, so $\alpha^*$ satisfies every $E_i$. By (C2), $\alpha^*$
satisfies some top-level cube. By (C2$'$), whenever $\alpha^*$ satisfies a node $P$ that was split, it satisfies some child of $P$ (in the
fallback case, use that $\alpha^*$ satisfies $B$). Following the tree down, $\alpha^*$ satisfies some leaf $Q_j$, so it satisfies $F_j$. This contradicts the unsatisfiability of $F_j$ established in (C3).
\end{proof}

\noindent\emph{Proof checking.} In leg~A each $F_i$ is solved by CaDiCaL~3.0.1~\cite{CaDiCaL2} with LRAT output, and the proof is
checked by \texttt{cake\_lpr}~\cite{cakelpr}. This checker is verified in CakeML down to machine code. In leg~B each $F_i$ is solved by
Kissat~4.0.4~\cite{Kissat24} with DRAT output and checked by \texttt{drat-trim}~\cite{DRATtrim}. This way the two legs share neither a
solver nor a proof checker. Before re-splitting was added, \texttt{drat-trim} needed up to $9.3$ hours on the proof of a single
oversized cube. With the caps above, oversized cubes are split before their proofs grow that large. The coverage formula (C2) is checked
in both ways in both legs. A separate audit script re-checks the bookkeeping: SHA-256 hashes of $B$ and the cube file are unchanged
since the start, there is exactly one certified record per cube, the literals in each record equal the cube in the file, and all
verdicts are present. For every split, the audit checks that each child's literals contain the parent's literals, and it re-proves the
split's coverage formula itself (Kissat and \texttt{drat-trim}) from the stored subcube files. It then checks that every top-level cube is
certified, either directly or recursively through its children. Proof files are deleted after a successful check; their sizes and SHA-256 hashes are kept in the logs.

\section{Validation on known values}\label{sec:valid}
We ran the pipeline on known values and reproduced each of them:

\begin{center}
\begin{tabular}{lll}
\toprule
value & result & checks \\
\midrule
$R(C_3,C_3,C_3)=17$ \cite{Rad} & $K_{16}$ SAT, $K_{17}$ UNSAT & \texttt{cake\_lpr}, \texttt{lrat-check}, \texttt{drat-trim} \\
$R(C_4,C_4,C_4)=11$ \cite{Rad} & $K_{10}$ SAT, $K_{11}$ UNSAT & \texttt{cake\_lpr}, \texttt{lrat-check}, \texttt{drat-trim} \\
$R(C_4,C_4,C_7)=12$ \cite{SunYLZ2,Tse3} & $K_{11}$ SAT, $K_{12}$ UNSAT & full CNF (both checkers); lazy C\&C legs A and B \\
$R(C_4,C_6,C_7)=13$ \cite{SunYWLX,Tse3} & $K_{12}$ SAT, $K_{13}$ UNSAT & full CNF (\texttt{cake\_lpr}, \texttt{lrat-check}, \texttt{drat-trim}); lazy C\&C leg A \\
\bottomrule
\end{tabular}
\end{center}
\noindent In addition, both lazy solvers found colourings of $K_{14}$ for $(C_4,C_7,C_7)$, and the C\&C driver correctly stopped
with a good colouring when run on $K_{12}$ for $(C_4,C_6,C_7)$. All such colourings passed both checkers. The clause checkers reject
deliberately corrupted clauses: a $C_4$-clause in a $C_7$-colour, a $C_7$-clause in the $C_4$-colour, a truncated cycle, and a
colouring clause offered as a cycle clause. Leg~B certified $R(C_4,C_4,C_7)\le 12$ ($K_{12}$, 121 cubes, all UNSAT, DRAT proofs verified by
\texttt{drat-trim}, coverage verified by \texttt{drat-trim} and \texttt{cake\_lpr}, audit passed). The cube-tree logic was tested on
$K_{12}$ by artificially capping solves at $2$\,s: this forced $13$ splits and $144$ leaves, and the audit passed. Deliberately deleting a
child from a split, or a leaf record, makes the audit fail.

\section{Results}\label{sec:results}

\begin{proof}[Proof of Theorem~\ref{thm:main}]
The lower bound is Proposition~\ref{prop:lb4}. For the upper bound we applied Theorem~\ref{thm:cert} with $n=15$:
\begin{itemize}
\item \emph{Leg A:} 988 cubes. All 988 are UNSAT, with $E_i$ checked and LRAT proofs verified by \texttt{cake\_lpr}. The coverage
certificate is verified by \texttt{cake\_lpr} and by \texttt{drat-trim}, and the audit passes. Total: 1672\,s lazy solving, 756\,s
offline proof generation, and 338\,s proof checking; 16 minutes wall-clock on 3 cores (Intel Xeon).
\item \emph{Leg B:} \TODO{1022 cubes; fill in counts, times, audit result when the run finishes.}
\end{itemize}
\end{proof}

\TODO{Theorem~\ref{thm:c6} (K21, leg A running since 18:45 PT, Oct 8, 2026; 1016 cubes): fill in when certified.
Estimated $\approx 30$ core-hours.}

\begin{table}[ht]
\caption{The four previously open entries of Table~XII of~\cite{Rad}.}\label{tab:results}
\begin{tabular}{lccl}
\toprule
 & lower bound & upper bound & status \\
\midrule
$R(C_3,C_7,C_7)$ & $25$ (Prop.~\ref{prop:odd}) & -- & open \\
$R(C_4,C_7,C_7)$ & $15$ (Prop.~\ref{prop:lb4}) & $15$ (certified) & $=15$ \\
$R(C_5,C_7,C_7)$ & $25$ (Prop.~\ref{prop:odd}) & -- & open \\
$R(C_6,C_7,C_7)$ & $21$ (Prop.~\ref{prop:lb6}) & \TODO{pending} & \TODO{pending} \\
\bottomrule
\end{tabular}
\end{table}

\begin{remark}
For $m\in\{3,5\}$ the corresponding computation on $K_{25}$ is beyond the methods used here. One possible route is a hybrid of the
Faudree--Schelten--Schiermeyer argument~\cite{FSS3} and SAT-checked local lemmas; we leave this open.
\end{remark}

\section{Reproducibility}\label{sec:repro}
The artifact \TODO{archive DOI} contains all generators, checkers, solver drivers, cube files, logs, per-cube status records with
proof hashes, the colourings, and a README with exact commands and SHA-256 checksums. Tool versions: CaDiCaL~3.0.1 (git
\texttt{c607304}), Kissat~4.0.4 (\texttt{8af8e56}), \texttt{cake\_lpr} (\texttt{a36874a}), \texttt{drat-trim} (\texttt{2e3b2dc}),
CnC/\texttt{march\_cu} (\texttt{705b60c}), PySAT~1.9.dev15, NetworkX~3.7, Python~3. The $K_{15}$ certificate can be re-created on
one machine in under an hour per leg. \TODO{update for leg B / K21.}

\section*{Use of AI tools}
\TODO{Disclosure placeholder. Describe how AI agents were used (literature search, code, running computations, drafting), what the
author checked personally, and confirm that the author takes full responsibility for the content. Follow the target journal's policy.}

\section*{Acknowledgements}
\TODO{}

\begin{thebibliography}{99}
\bibitem{CaDiCaL2} A.~Biere, T.~Faller, K.~Fazekas, M.~Fleury, N.~Froleyks and F.~Pollitt, CaDiCaL 2.0, in: \emph{Computer Aided
Verification (CAV 2024)}, LNCS 14682, Springer, 2024, 133--152.
\bibitem{Kissat24} A.~Biere, T.~Faller, K.~Fazekas, M.~Fleury, N.~Froleyks and F.~Pollitt, CaDiCaL, Gimsatul, IsaSAT and Kissat entering
the SAT Competition 2024, in: \emph{Proc.\ SAT Competition 2024}, Dept.\ of Computer Science Report Series B, vol.\ B-2024-1,
University of Helsinki, 2024, 8--10.
\bibitem{CGLR} J.~Crawford, M.~Ginsberg, E.~Luks and A.~Roy, Symmetry-breaking predicates for search problems, in: \emph{Proc.\ KR'96},
Morgan Kaufmann, 1996, 148--159.
\bibitem{EFRS1} P.~Erd\H{o}s, R.\,J.~Faudree, C.\,C.~Rousseau and R.\,H.~Schelp, Generalized Ramsey theory for multiple colors,
\emph{J.\ Combin.\ Theory Ser.\ B} \textbf{20} (1976), 250--264.
\bibitem{FSS3} R.\,J.~Faudree, A.~Schelten and I.~Schiermeyer, The Ramsey number $r(C_7,C_7,C_7)$, \emph{Discuss.\ Math.\ Graph
Theory} \textbf{23} (2003), 141--158.
\bibitem{IPASIRUP} K.~Fazekas, A.~Niemetz, M.~Preiner, M.~Kirchweger, S.~Szeider and A.~Biere, IPASIR-UP: user propagators for CDCL,
in: \emph{SAT 2023}, LIPIcs 271, 2023, 8:1--8:13.
\bibitem{HKWB} M.\,J.\,H.~Heule, O.~Kullmann, S.~Wieringa and A.~Biere, Cube and conquer: guiding CDCL SAT solvers by lookaheads,
in: \emph{HVC 2011}, LNCS 7261, Springer, 2012, 50--65.
\bibitem{KoSS2} Y.~Kohayakawa, M.~Simonovits and J.~Skokan, The 3-colored Ramsey number of odd cycles, \emph{J.\ Combin.\ Theory Ser.\ B}
(as cited in~\cite{Rad}). \TODO{verify volume/pages}
\bibitem{LidP} B.~Lidick\'y and F.~Pfender, Semidefinite programming and Ramsey numbers, \emph{SIAM J.\ Discrete Math.} \textbf{35}(4)
(2021); arXiv:1704.03592. \TODO{pages}
\bibitem{Rad} S.\,P.~Radziszowski, Small Ramsey numbers, \emph{Electron.\ J.\ Combin.}, Dynamic Survey DS1, revision~18 (April 24,
2026), \url{https://doi.org/10.37236/21}.
\bibitem{SunYWLX} Sun Yongqi, Yang Yuansheng, Wang Wei, Li Bingxi and Xu Feng, Study of three color Ramsey numbers
$R(C_{m_1},C_{m_2},C_{m_3})$ (in Chinese), \emph{J.\ Dalian Univ.\ Technol.} \textbf{46} (2006), 428--433.
\bibitem{SunYLZ2} Sun Yongqi, Yang Yuansheng, Lin Xiaohui and Zheng Wenping, On the three color Ramsey numbers $R(C_m,C_4,C_4)$,
\emph{Ars Combin.} \textbf{84} (2007), 3--11.
\bibitem{cakelpr} Y.\,K.~Tan, M.\,J.\,H.~Heule and M.\,O.~Myreen, Verified propagation redundancy and compositional UNSAT checking in
CakeML, \emph{Int.\ J.\ Softw.\ Tools Technol.\ Transf.} \textbf{25} (2023), 167--184.
\bibitem{Tse3} K.-K.~Tse, A note on some Ramsey numbers $R(C_p,C_q,C_r)$, \emph{J.\ Combin.\ Math.\ Combin.\ Comput.} \textbf{62}
(2007), 189--192.
\bibitem{We3} W.\,J.~Wesley, New bounds for some small multicolor Ramsey numbers, arXiv:2509.03784, 2025.
\bibitem{DRATtrim} N.~Wetzler, M.\,J.\,H.~Heule and W.\,A.~Hunt Jr., DRAT-trim: efficient checking and trimming using expressive
clausal proofs, in: \emph{SAT 2014}, LNCS 8561, Springer, 2014, 422--429.
\end{thebibliography}

\end{document}
